\begin{thebibliography}{25}

\bibitem[1]{ref1}
Anurag Anshu, Dmitry Gavinsky, Rahul Jain, Srijita Kundu, Troy Lee, Priyanka Mukhopad- hyay, Miklos Santha, and Swagato Sanyal. A composition theorem for randomized query complexity. In Proceedings of the 37th Conference on Foundations of Software Technology and Theoretical Computer Science (FSTTCS), pages 10:1–10:13. Schloss Dagstuhl, 2017. doi:10.4230/LIPIcs.FSTTCS.2017.10.

\bibitem[2]{ref2}
Andrew Bassilakis, Andrew Drucker, Mika Göös, Lunjia Hu, Weiyun Ma, and Li-Yang Tan. The power of many samples in query complexity. In Proceedings of the 47th International Colloquium on Automata, Languages, and Programming (ICALP), volume 168, pages 9:1–9:18. Schloss Dagstuhl, 2020. doi:10.4230/LIPIcs.ICALP.2020.9.

\bibitem[3]{ref3}
Shalev Ben-David and Eric Blais. A new minimax theorem for randomized algorithms. In Proceedings of the 61st Symposium on Foundations of Computer Science (FOCS), pages 403–411, 2020. doi:10.1109/FOCS46700.2020.00045.

\bibitem[4]{ref4}
Shalev Ben-David and Eric Blais. A tight composition theorem for the randomized query complexity of partial functions. In Proceedings of the 61st Symposium on Foundations of Computer Science (FOCS), pages 240–246, 2020. doi:10.1109/FOCS46700.2020.00031.

\bibitem[5]{ref5}
Shalev Ben-David, Mika Göös, Robin Kothari, and Thomas Watson. When is amplification necessary for composition in randomized query complexity? In Proceedings of the 22nd International Conference on Randomization and Computation (RANDOM), volume 176, pages 28:1–28:16. Schloss Dagstuhl, 2020. doi:10.4230/LIPIcs.APPROX/RANDOM.2020.28.

\bibitem[6]{ref6}
Shalev Ben-David and Robin Kothari. Randomized query complexity of sabotaged and com- posed functions. Theory of Computing, 14(1):1–27, 2018. doi:10.4086/toc.2018.v014a005.

\bibitem[7]{ref7}
Eric Blais and Joshua Brody. Optimal separation and strong direct sum for randomized query complexity. In Proceedings of the 34th Computational Complexity Conference (CCC), pages 29:1–29:17. Schloss Dagstuhl, 2019. doi:10.4230/LIPIcs.CCC.2019.29.

\bibitem[8]{ref8}
Joshua Brody, Jae Tak Kim, Peem Lerdputtipongporn, and Hariharan Srinivasulu. A strong XOR lemma for randomized query complexity. Technical report, arXiv, 2020. arXiv:2007. 05580.

\bibitem[9]{ref9}
Harry Buhrman and Ronald de Wolf. Complexity measures and decision tree complexity: A survey. Theoretical Computer Science, 288(1):21–43, 2002. doi:10.1016/S0304-3975(01) 00144-X.

\bibitem[10]{ref10}
Chinmoy Dutta and Jaikumar Radhakrishnan. Lower bounds for noisy wireless networks using sampling algorithms. In Proceedings of the 49th Symposium on Foundations of Computer Science (FOCS), pages 394–402. IEEE, 2008. doi:10.1109/FOCS.2008.72.

\bibitem[11]{ref11}
William Evans and Nicholas Pippenger. Average-case lower bounds for noisy boolean decision trees. SIAM Journal on Computing, 28(2):433–446, 1998. doi:10.1137/S0097539796310102.

\bibitem[12]{ref12}
Uriel Feige, Prabhakar Raghavan, David Peleg, and Eli Upfal. Computing with noisy informa- tion. SIAM Journal on Computing, 23(5):1001–1018, 1994. doi:10.1137/S0097539791195877.

\bibitem[13]{ref13}
Dmitry Gavinsky, Troy Lee, Miklos Santha, and Swagato Sanyal. A composition theorem for randomized query complexity via max-conflict complexity. In Proceedings of the 46th International Colloquium on Automata, Languages, and Programming (ICALP), pages 64:1– 64:13. Schloss Dagstuhl, 2019. doi:10.4230/LIPIcs.ICALP.2019.64.

\bibitem[14]{ref14}
Mika Göös, Shachar Lovett, Raghu Meka, Thomas Watson, and David Zuckerman. Rectangles are nonnegative juntas. SIAM Journal on Computing, 45(5):1835–1869, 2016. doi:10.1137/ 15M103145X.

\bibitem[15]{ref15}
Navin Goyal and Michael Saks. Rounds vs. queries tradeoff in noisy computation. Theory of Computing, 6(1):113–134, 2010. doi:10.4086/toc.2010.v006a006.

\bibitem[16]{ref16}
Rahul Jain and Hartmut Klauck. The partition bound for classical communication complexity and query complexity. In Proceedings of the 25th Conference on Computational Complexity (CCC), pages 247–258. IEEE, 2010. doi:10.1109/CCC.2010.31.

\bibitem[17]{ref17}
Rahul Jain, Hartmut Klauck, and Miklos Santha. Optimal direct sum results for deterministic and randomized decision tree complexity. Information Processing Letters, 110(20):893–897, 2010. doi:10.1016/j.ipl.2010.07.020.

\bibitem[18]{ref18}
Stasys Jukna. Boolean Function Complexity: Advances and Frontiers, volume 27 of Algorithms and Combinatorics. Springer, 2012.

\bibitem[19]{ref19}
Jedrzej Kaniewski, Troy Lee, and Ronald de Wolf. Query complexity in expectation. In Proceedings of the 42nd International Colloquium on Automata, Languages, and Programming (ICALP), pages 761–772. Springer, 2015. doi:10.1007/978-3-662-47672-7\_62.

\bibitem[20]{ref20}
Claire Kenyon and Valerie King. On boolean decision trees with faulty nodes. Random Structures and Algorithms, 5(3):453–464, 1994. doi:10.1002/rsa.3240050306.

\bibitem[21]{ref21}
Troy Lee, Rajat Mittal, Ben Reichardt, Robert Špalek, and Mario Szegedy. Quantum query complexity of state conversion. In Proceedings of the 52nd Symposium on Foundations of Computer Science (FOCS), pages 344–353. IEEE, 2011. doi:10.1109/FOCS.2011.75.

\bibitem[22]{ref22}
Ilan Newman. Computing in fault tolerant broadcast networks and noisy decision trees. Random Structures and Algorithms, 34(4):478–501, 2009. doi:10.1002/rsa.20240.

\bibitem[23]{ref23}
Ben Reichardt. Reflections for quantum query algorithms. In Proceedings of the 22nd Symposium on Discrete Algorithms (SODA), pages 560–569. SIAM, 2011.

\bibitem[24]{ref24}
Petr Savický. On determinism versus unambiquous nondeterminism for decision trees. Technical Report TR02-009, Electronic Colloquium on Computational Complexity (ECCC), 2002. URL: http://eccc.hpi-web.de/report/2002/009/.

\bibitem[25]{ref25}
Avishay Tal. Properties and applications of boolean function composition. In Proceedings of the 4th Conference on Innovations in Theoretical Computer Science (ITCS), pages 441–454. ACM, 2013. doi:10.1145/2422436.2422485.
\end{thebibliography}
